\chapter{Introduction to CS 135 and Racket}

\section{What this course is about}

CS 135 is the first course in the CS major sequence at Waterloo. Rather than
starting with an industrial language such as C++ or Java, it uses
\textbf{Racket}, a dialect of Scheme/Lisp, through a series of restricted
\emph{teaching languages} in the DrRacket environment:
\begin{itemize}
  \item \textbf{Beginning Student} (BSL) --- functions, constants, \code{cond}, structures, lists.
  \item \textbf{Beginning Student with List Abbreviations} --- adds \code{list} and quote notation.
  \item \textbf{Intermediate Student} --- adds \code{local}.
  \item \textbf{Intermediate Student with lambda} (ISL+) --- adds anonymous functions and functions as first-class values.
\end{itemize}

The emphasis is on \emph{design}: how to go systematically from a problem
statement to a correct, readable, tested program. The tool for this is the
\textbf{design recipe} (Chapter~\ref{ch:design-recipe}). The programming style
is \textbf{functional}: there is no mutation and no assignment statement. A
program is a collection of function definitions, and running it means
\emph{evaluating an expression}.

\begin{aside}
Why a functional language first? Because without mutation, the meaning of a
program can be explained entirely by algebraic substitution, the same way you
simplify $f(3)$ in high-school algebra. That makes reasoning about correctness
much easier, and everyone starts on equal footing regardless of prior
programming experience.
\end{aside}

\section{From mathematical functions to Racket}

In mathematics we write $f(x) = x^2$ and $g(x,y) = x + y$, and evaluate
$g(f(2), f(3)) = g(4, 9) = 13$. Racket uses \textbf{prefix notation} with
parentheses around every application:

\begin{lstlisting}[style=racket]
(define (f x) (* x x))
(define (g x y) (+ x y))
(g (f 2) (f 3))   ; => 13
\end{lstlisting}

Points to note:
\begin{itemize}
  \item The function (or operator) always comes first: \code{(+ 1 2)}, not \code{1 + 2}.
  \item Every parenthesis matters. \code{(f 2)} is an application; \code{((f 2))} tries to apply the number 4 as a function and is an error.
  \item There is no operator precedence to memorize: \code{(* 2 (+ 3 4))} is unambiguous.
  \item Many built-ins are variadic: \code{(+ 1 2 3 4)} is 10.
\end{itemize}

\subsection{Numbers}
Racket numbers are \emph{exact} whenever possible. Integers are unbounded and
rational arithmetic is exact: \code{(/ 1 3)} is the rational $1/3$, not
$0.333\ldots$. Only operations like \code{sqrt} on non-perfect squares produce
\emph{inexact} numbers, displayed with a \code{\#i} prefix (e.g.\
\code{\#i1.414...}). Inexact numbers cannot be compared reliably with
\code{check-expect}; use \code{check-within} instead.

Useful built-ins: \code{quotient}, \code{remainder}, \code{modulo},
\code{expt}, \code{sqrt}, \code{abs}, \code{max}, \code{min}, \code{floor},
\code{ceiling}, \code{exact->inexact}, \code{even?}, \code{odd?}, \code{zero?}.

\section{Expressions, constants and functions}

A Racket program is a sequence of \textbf{definitions} and \textbf{expressions}.

\begin{definition}[Constant definition]
\code{(define name expr)} evaluates \code{expr} to a value and binds \code{name} to it.
\end{definition}

\begin{definition}[Function definition]
\code{(define (fname x1 ... xn) body)} creates a function with
\textbf{parameters} \code{x1} \ldots\ \code{xn}. When applied to
\textbf{arguments}, the body is evaluated with each parameter replaced by the
corresponding argument value.
\end{definition}

\begin{lstlisting}[style=racket]
(define tax-rate 13/100)
(define (price-with-tax p)
  (+ p (* p tax-rate)))
(price-with-tax 100)   ; => 113
\end{lstlisting}

Constants make programs easier to read and change: if the tax rate changes, it
is updated in one place. Avoid ``magic numbers'' in function bodies.

\section{Syntax and semantics}

Every language has a \textbf{syntax} (which programs are well-formed) and a
\textbf{semantics} (what a well-formed program means). For natural languages:
``Trombones fly hungrily'' is syntactically fine but semantically nonsense.

\subsection{Syntax}
Racket syntax is very small. Ignoring later additions, an expression is one of:
\begin{itemize}
  \item a value (number, boolean, string, symbol, \ldots);
  \item an identifier (name of a constant or parameter);
  \item an application \code{(f e1 ... en)} where \code{f} is a function name and each \code{ei} is an expression;
  \item a special form such as \code{(cond [q1 a1] ...)}, \code{(and ...)}, \code{(or ...)}.
\end{itemize}
Errors such as unbalanced parentheses or \code{(define 5 x)} are
\emph{syntax errors}: DrRacket refuses to run the program at all.

\subsection{Semantics: the substitution model}
CS 135 gives Racket a precise semantics as a set of \textbf{rewrite (reduction)
rules}. A program is evaluated by repeatedly finding the leftmost innermost
subexpression that can be simplified, rewriting it, and continuing until only
values remain. Each rewrite is one \emph{step}, written $\Rightarrow$.

\begin{important}[Substitution rules (Beginning Student)]
\begin{enumerate}
  \item \textbf{Built-in application}: \code{(f v1 ... vn)} $\Rightarrow$ \code{v}, where \code{f} is built in, the \code{vi} are values, and \code{v} is the result of applying \code{f}.
  \item \textbf{User-defined application}: if \code{(define (f x1 ... xn) exp)} occurs earlier, then \code{(f v1 ... vn)} $\Rightarrow$ \code{exp'}, where \code{exp'} is \code{exp} with every \code{xi} replaced by \code{vi}.
  \item \textbf{Constant}: if \code{(define id val)} occurs earlier, then \code{id} $\Rightarrow$ \code{val}.
  \item Arguments are evaluated left to right, and \emph{all arguments must be values} before rules 1 or 2 apply.
\end{enumerate}
\end{important}

Constant definitions are themselves evaluated: \code{(define x (+ 1 2))}
becomes \code{(define x 3)} before later expressions can use \code{x}.

\begin{example}[Stepper trace]
Using the definitions of \code{f} and \code{g} above:
\begin{lstlisting}[style=racket]
(g (f 2) (f 3))
=> (g (* 2 2) (f 3))     ; rule 2 on (f 2)
=> (g 4 (f 3))           ; rule 1
=> (g 4 (* 3 3))         ; rule 2 on (f 3)
=> (g 4 9)               ; rule 1
=> (+ 4 9)               ; rule 2 on g
=> 13                    ; rule 1
\end{lstlisting}
\end{example}

It helps to picture the expression as a tree: each application is a node whose
children are its arguments. Evaluation works from the leaves up, finishing
the left argument before starting the right one.

\begin{center}
\begin{tikzpicture}[level distance=13mm,
  level 1/.style={sibling distance=40mm},
  app/.style={draw,rounded corners,fill=blue!6,font=\ttfamily\small,minimum height=6mm,minimum width=8mm},
  lf/.style={draw,circle,font=\ttfamily\small,minimum size=6mm,inner sep=1pt},
  val/.style={font=\scriptsize,text=blue!50!black,align=left}]
  \node[app] (g) {(g \_ \_)}
    child {node[app] (f1) {(f \_)} child {node[lf] {2}}}
    child {node[app] (f2) {(f \_)} child {node[lf] {3}}};
  \node[val,left=2mm of f1] {steps 1--2\\$\Rightarrow$ \texttt{4}};
  \node[val,right=2mm of f2] {steps 3--4\\$\Rightarrow$ \texttt{9}};
  \node[val,right=2mm of g] {steps 5--6\\$\Rightarrow$ \texttt{13}};
\end{tikzpicture}
\end{center}

Note that \code{(f 2)} is rewritten to \code{(* 2 2)} in one step (the body
with substitution), not directly to \code{4}. Exam questions frequently ask for
the \emph{exact} sequence of steps, so be precise about this. DrRacket's
\textbf{Stepper} tool shows exactly these steps and is the best way to check
your traces.

\begin{example}[Trace with a constant]
\begin{lstlisting}[style=racket]
(define tax-rate 13/100)
(price-with-tax 200)
=> (+ 200 (* 200 tax-rate))
=> (+ 200 (* 200 13/100))
=> (+ 200 26)
=> 226
\end{lstlisting}
\end{example}

\subsection{Run-time errors}
A syntactically valid program can still fail during evaluation, e.g.\
\code{(/ 1 0)} or \code{(f 1 2)} when \code{f} takes one argument. In the
substitution model this shows up as an expression that is not a value but to
which no rule applies; evaluation stops with an error.

\section{Identifiers and scope (preview)}
Parameter names are local to their function: the \code{x} in \code{f} has
nothing to do with an \code{x} in \code{g}. Constants defined at the top level
are visible everywhere after their definition. A parameter with the same name
as a constant \emph{shadows} the constant inside that function body. Chapter~\ref{ch:local}
makes this precise.

Identifiers in Racket may contain almost any character except spaces,
parentheses and a few others, so \code{price-with-tax}, \code{empty?},
\code{->string} and \code{x*} are all legal names. Conventionally, names are
lowercase words separated by hyphens, and predicates end in \code{?}.

\section{Programs in DrRacket}
\begin{itemize}
  \item The \textbf{definitions window} holds the program; \textbf{Run} evaluates it top to bottom and shows results of top-level expressions in the \textbf{interactions window}.
  \item Tests (\code{check-expect}) are collected and run after the whole program, so a test may appear before the function it tests.
  \item Choose the correct language level from \emph{Language $\to$ Choose Language}. Using a feature not in the current level is a syntax error; this is intentional, to catch typical beginner mistakes with clear messages.
\end{itemize}

\section{Common pitfalls}
\begin{itemize}
  \item Extra or missing parentheses: \code{(define (f x) (x))} applies \code{x} as a function.
  \item Infix habits: \code{(2 + 3)} is an error.
  \item Forgetting that all arguments are evaluated first, so a trace that rewrites the outer function too early is wrong.
  \item Comparing inexact numbers with \code{check-expect}.
\end{itemize}

\section{Summary}
\begin{itemize}
  \item CS 135 teaches program design using Racket teaching languages, without mutation.
  \item Racket uses fully parenthesized prefix notation; a program is definitions plus expressions.
  \item Semantics is given by the substitution model: evaluate arguments left to right to values, then substitute into the function body; constants are replaced by their values.
  \item Know how to produce a precise step-by-step trace; the Stepper is the reference.
\end{itemize}
